% =========================================================
% The Two Questions Before You Write Anything
% =========================================================
\paragraph{The Two Questions Before You Write Anything.}

Before writing a single symbol, ask two questions in order.

\begin{tcolorbox}[colback=propbox, colframe=propborder, arc=2pt,
  left=6pt, right=6pt, top=4pt, bottom=4pt]
\textbf{Question 1.} \emph{What is the logical form of the statement?}

Read the claim and identify its outermost logical structure:
universal, conditional, biconditional, existential, uniqueness,
equality, or structural assertion.

\medskip
\textbf{Question 2.} \emph{What proof strategy does this form demand?}

The form determines the skeleton of the proof almost mechanically,
before any mathematics is done.
\end{tcolorbox}

\begin{remark}[Why order matters --- and why most students get it backwards]
Question~2 cannot be answered before Question~1, but almost every
beginner tries to answer it first. The default instinct is to reach for
a familiar technique --- proof by contradiction is the usual reflex ---
and start writing. This is the wrong order.

Reading the statement \emph{is} the first mathematical act of the proof.
The logical form of the statement is not preliminary bookkeeping; it
is the information that determines the entire structure of the argument.
A proof whose structure mismatches the statement's form is not just
inelegant --- it is often invalid or at best needlessly circular.

Until you are fluent enough to read the form automatically, slow down
deliberately before every proof. Write the statement out. Identify its
outermost connective. Then and only then choose a strategy.
\end{remark}

\begin{example}[Applying the two questions]
\textbf{Statement.} \emph{The identity element of a group is unique.}

\medskip
\noindent\textbf{Q1. Logical form.}
Strip the natural-language phrasing down to its logical skeleton.
``The identity is unique'' means: for every group $G$ and for any two
elements $e, e' \in G$ that both satisfy the identity axiom, we have
$e = e'$.
The outermost structure is a universal statement with a uniqueness
conclusion: $\forall G,\; \forall e, e' \in G$,
[both satisfy identity axiom] $\Rightarrow e = e'$.

\medskip
\noindent\textbf{Q2. Strategy.}
A universally quantified uniqueness claim demands the
\emph{assume-two-and-compare} pattern: introduce an arbitrary group~$G$,
assume two objects $e$ and $e'$ both satisfying the definition, and
derive $e = e'$.

The strategy is forced entirely by the form. No creative insight is
required to select it. The selection is a consequence of reading.
\end{example}

\begin{remark}[What ``logical form'' means in practice]
You do not need to write out the full first-order formula every time.
You need to identify one thing: the \emph{outermost} structure.

Is the statement ``for all $x$, \ldots''? Then introduce an arbitrary $x$.
Is it ``$P$ implies $Q$''? Then assume $P$.
Is it ``there exists $x$''? Then construct a witness.
Is it ``there exists a unique $x$''? Then split into existence and
uniqueness.
Is it ``$A = B$ as sets''? Then prove double inclusion.
Is it ``$a = b$ as elements in an algebraic structure''? Then either do
algebra or use satisfy-and-cite.

The statement map in the next section makes this correspondence explicit.
\end{remark}
